\documentclass[11pt]{article}
\usepackage[a4paper,margin=1in]{geometry}
\usepackage{amsmath,amssymb,amsthm,booktabs,microtype}
\newtheorem{theorem}{Theorem}[section]
\newtheorem{corollary}[theorem]{Corollary}
\theoremstyle{remark}\newtheorem{remark}[theorem]{Remark}
\title{Canonical Modulus-Complete Heads for Projection-Free Spectral Factors}
\author{Bin Wang}\date{July 2026}
\begin{document}\maketitle
\begin{abstract}
The rank-growing tail theorem of RH-282 requires an actual noisy spectral
head, not a fitted root list.  For any Hilbert--Schmidt operator and threshold
$q>0$, we define the head to contain every algebraic eigenvalue of modulus
strictly larger than $q$.  The head is finite, conjugation complete for real
operators, and has cardinality at most the squared spectral mass divided by
$q^2$.  It is the unique smallest spectral submultiset whose complement has
spectral radius at most $q$.  The strict threshold convention also resolves
ties.  Its genus-one product is an exact finite factor of $\det_2$, independent
of eigenvector conditioning.  Canonicality is relative to $q$ and does not
identify the selected roots with the deterministic monodromy shell.
\end{abstract}

\section{Definition}
Let $A\in\mathcal S_2(H)$ and let $(\mu_j)$ be its nonzero eigenvalues with
algebraic multiplicity.  For $q>0$ set
\[
 H_q(A)=\{\mu_j:|\mu_j|>q\},\qquad
 T_q(A)=\{\mu_j:|\mu_j|\le q\}.
\]
The equality case belongs to the tail.  This strict convention is part of the
definition.

\section{Minimal head theorem}
\begin{theorem}\label{thm:head}
Write $M=\sum_j|\mu_j|^2$.  Then $H_q(A)$ is finite and
\[
 \#H_q(A)\le M/q^2\le\|A\|_2^2/q^2.
\]
If a finite spectral submultiset $F$ has the property that every eigenvalue
outside $F$ has modulus at most $q$, then $H_q(A)\subseteq F$.  Thus
$H_q(A)$ is the unique smallest admissible head.
\end{theorem}
\begin{proof}
Every member of $H_q$ contributes more than $q^2$ to the convergent squared
sum, proving finiteness and the count.  If $F$ omitted an eigenvalue with
modulus greater than $q$, its complement would violate the radius condition.
Hence every admissible $F$ contains $H_q$.
\end{proof}

\begin{corollary}[real symmetry]
If $A$ commutes with a real conjugation, then $H_q(A)$ is closed under complex
conjugation, with equal algebraic multiplicities.  Its finite canonical
product has real Taylor coefficients.
\end{corollary}
\begin{proof}
The algebraic spectrum of a real operator is conjugation invariant, and
modulus is unchanged by conjugation.
\end{proof}

\section{Exact determinant split}
The genus-one products
\[
 D_{H,q}(z)=\prod_{\mu\in H_q}(1-z\mu)e^{z\mu},\qquad
 D_{T,q}(z)=\prod_{\mu\in T_q}(1-z\mu)e^{z\mu}
\]
are well defined, the first finite and the second locally normally
convergent.  RH-234 gives the exact identity
\[
 \det_2(I-zA)=D_{H,q}(z)D_{T,q}(z).
\]
No invariant-subspace basis is needed.

\begin{remark}
The threshold $q$ is an explicit design parameter.  The theorem is canonical
after $q$ is fixed, not a claim that the dynamics chooses $q=1/2$.  It also
does not prove that $H_q$ equals the RH-272 counterloop atoms.  Gates A--E
remain false/open.
\end{remark}
\end{document}
